\section{Uso de \textit{ik / jij / hij / wij / jullie / zij}}
\textbf{Enfoque:} Combinar sujeto correcto con forma verbal y orden V2.

\textbf{Tips}
\begin{itemize}[leftmargin=*]
  \item V2: el verbo conjugado va en segunda posición.\ \textit{Vandaag werk ik thuis. Morgen speel jij voetbal.}
  \item \textit{je/jij}: \textit{jij} para énfasis; \textit{je} neutro.\ \textit{u} formal (usa la forma con \textit{-t}).
  \item Plural: \textit{wij/jullie/zij} usan infinitivo.
\end{itemize}

\textbf{Práctica}
\begin{itemize}[leftmargin=*]
  \item Rueda de sujetos: di una misma acción con las 6 personas.
  \item Reordena frases con adverbio inicial para respetar V2.
\end{itemize}
